% Source: https://artofproblemsolving.com/community/c3h2725367_infinite_continued_fraction_general_case
% Stable identity: aops-topic-2725367-q01
The tail after the first $k$ is the same continued fraction, so $x=k+\dfrac1x$, i.e. $x^2-kx-1=0$. Hence $x=\dfrac{k\pm\sqrt{k^2+4}}{2}$. Since $x\ge k>0$ and $\sqrt{k^2+4}>k$, the minus sign gives a negative number, so
\[x=\boxed{\dfrac{k+\sqrt{k^2+4}}{2}}.\]
(The convergents $x_1=k$, $x_{n+1}=k+\dfrac1{x_n}$ do converge to this value: the map $t\mapsto k+\dfrac1t$ is a contraction on $t\ge k$ when $k>1$, and in general the even and odd convergents are monotone and squeeze to the same limit.)
